% Slide: Critical Thinking Questions and Answers
\begin{frame}{Critical Thinking Questions and Answers}
  \textbf{Question 1: Can Canada remain competitive in a rapidly changing global economy? Why or why not?}
  \pause
  \textbf{Answer:}
  \begin{itemize}
    \item \textbf{Strengths:}
      \begin{itemize}
        \item Abundance of natural resources (e.g., oil, natural gas, minerals).
        \item Skilled and well-educated workforce.
        \item Stable political and financial systems attracting foreign investment.
      \end{itemize}
    \item \textbf{Challenges:}
      \begin{itemize}
        \item Global competition from emerging economies like China and India.
        \item Rapid technological disruption requiring adaptation.
        \item Need for environmental sustainability and reduced reliance on fossil fuels.
      \end{itemize}
    \item \textbf{Policy Recommendations:}
      \begin{itemize}
        \item Invest in renewable energy and green technologies.
        \item Encourage innovation and R&D in AI, biotechnology, and advanced manufacturing.
        \item Strengthen trade relationships and diversify exports.
      \end{itemize}
  \end{itemize}
\end{frame}

\begin{frame}{Critical Thinking Questions and Answers (cont'd)}
  \textbf{Question 2: To what extent should Canada rely on natural resources for economic growth?}
  \pause
  \textbf{Answer:}
  \begin{itemize}
    \item \textbf{Short-term reliance:}
      \begin{itemize}
        \item Natural resources drive economic stability, providing export revenues and jobs.
      \end{itemize}
    \item \textbf{Long-term risks:}
      \begin{itemize}
        \item Vulnerability to global commodity price fluctuations.
        \item Environmental concerns and global climate policy pressures.
      \end{itemize}
    \item \textbf{Path Forward:}
      \begin{itemize}
        \item Transition from resource-based growth to knowledge-based industries like technology and services.
        \item Invest in renewable energy and clean technologies.
        \item Support resource-dependent regions with diversification and reskilling policies.
      \end{itemize}
  \end{itemize}
\end{frame}

\begin{frame}{Critical Thinking Questions and Answers (cont'd)}
  \textbf{Question 3: What role does technology play in shaping future economic policies?}
  \pause
  \textbf{Answer:}
  \begin{itemize}
    \item \textbf{Opportunities:}
      \begin{itemize}
        \item Increases productivity through automation and AI.
        \item Creates new industries (e.g., green technology, fintech, and biotech).
        \item Enhances global connectivity for Canadian businesses.
      \end{itemize}
    \item \textbf{Challenges:}
      \begin{itemize}
        \item Job displacement in low-skilled industries.
        \item Regulation of technology, addressing concerns like privacy and cybersecurity.
        \item Bridging the digital divide for equitable access to technology.
      \end{itemize}
    \item \textbf{Policy Recommendations:}
      \begin{itemize}
        \item Invest in digital infrastructure and rural connectivity.
        \item Implement workforce reskilling programs to prepare workers for tech-driven industries.
        \item Collaborate with tech companies to regulate innovation ethically and sustainably.
      \end{itemize}
  \end{itemize}
\end{frame}
